% ===================================================================
%  اللوحة نمرة ٥ — البيت وبناه.
%
%  Traduction, faite en 2026, en arabe egyptien ecrit, sur l'ido de
%  texto/io. Regles de la colonne : voir 10-lawha-01.tex. Le tableau
%  est un DIALOGUE : on garde \textsc{} pour le nom du parleur,
%  comme la colonne arabe.
%
%  C'EST LE PENDANT EXACT DU TABLEAU 6 CANTONAIS, ET LE SUFFIXE -جي
%  Y FAIT CE QUE 佬 Y FAISAIT :
%
%    عربجي (135)   le charretier   — de عربية, la voiture
%    كهربجي (124)  l'electricien
%
%  et autour d'eux les noms de metier egyptiens en -ّاي / -ّان :
%  البنّا (108) le macon, النقّاش (102) le peintre, القزّاز (96) le
%  vitrier, السمكري (121) le zingueur, المبيّض (99) le platrier,
%  الجنايني (129) le jardinier, السايس (133) le palefrenier, الشيّال
%  (146) le manœuvre, الحفّار (81) le terrassier. Pas un n'est un
%  participe savant : ce sont des noms de metier faits sur l'ouvrage.
%
%  ET LE MAITRE D'ŒUVRE S'APPELLE الأسطى, le titre que l'Egypte donne
%  au maitre artisan. C'est ce qu'est monsieur Blanc dans le livret,
%  et le mot le dit mieux que « السيد ».
%
%  LES OUTILS PARLENT TURC ET ITALIEN, ET AUCUN NE SE DIT EN ARABE
%  STANDARD :
%
%    كوريك (82)    la pelle       du turc kürek
%    بوية (103)    la peinture    du turc boya
%    طرمبة (58)    la pompe
%    جردل (114)    le seau        MSA : دلو
%    زرديّة (128)   la pince
%    سندال (127)   l'enclume
%    ترباس (53)    le verrou
%    كالون (6)     la serrure
%    شكارة (71)    le sac
%    سقالة (110)   l'echafaudage
%    بلطة (116)    la hache
%    فارة (91)     le rabot
%    مسطرين (100)  la truelle
%    باركيه (16)   le parquet
%    شيش (19)      les persiennes
%    بلكونة (28)   le balcon      MSA : شرفة
%    درابزين (15)  la rampe
%    بدروم (3)     la cave        MSA : قبو
%    كارّو (136)    le chariot
%    مونة (65)     le mortier
%    طوب (61)      la brique
%
%  ET LA MAISON SE DIVISE EN أُوَض, NON EN غرف : أوضة النوم (g),
%  أوضة السفرة (c), أوضة العيال (k), أوضة الضيوف (i). C'est le mot
%  du tableau 1 — celui qui figure dans la liste de controle de la
%  colonne — et le tableau de la maison est l'endroit ou il revient
%  huit fois.
%
%  L'EAU EST مية, NON ماء. Le mot le plus ordinaire de la langue est
%  aussi celui qui separe le plus vite les deux registres.
% ===================================================================

%%K t05-apar-1 apar
\VUsaut{6.0mm}
\VUtitre{15.0pt}{60}{اللوحة نمرة ٥}
\VUsaut{1.4mm}
\VUtitre{11.4pt}{0}{البيت وبناه.}
\VUsaut{2.2mm}

%%K t05-tit-1 sub
\VUpk{10.2pt}{40}{مقابلة كويسة.}
\VUsaut{3.0mm}

%%K t05-01-1 p
\textsc{يوحنا}. --- يا عمّي، نفسي أوي أعرف \VUgras{البيت} بيتبني إزاي.

%%K t05-01-2 p
\textsc{العم} \textsuperscript{(80)}. --- ده سهل أوي يا حبيبي، تعالى معايا،
هوريك البيت اللي الأستاذ برناردوس \textsuperscript{(79)} بيبنيه واللي أنا هسكن فيه.

%%K t05-01-3 p
\textsc{يوحنا}. --- ومين اللي رسم \VUgras{رسم} \textsuperscript{(76)} البيت ده؟

%%K t05-01-4 p
\textsc{العم}. --- \VUgras{المهندس المعماري} \textsuperscript{(75)}؛ قدّم رسم واضح
أوي واتوافق عليه، و\VUgras{المقاول} \textsuperscript{(77)} بدأ الشغل.

%%K t05-01-5 p
\textsc{يوحنا}. --- ومين المقاول؟

%%K t05-01-6 p
\textsc{العم}. --- هو الأسطى بلانكو، والد صاحبك يعقوب. وهو بيدير العمال مع
المهندس روفو.

%%K t05-01-7 p
أهو! أهو الأسطى بلانكو جاي في وقته بـ\VUgras{مسطرته} \textsuperscript{(78)}، والمهندس
روفو بالرسم بتاعه.

%%K t05-01-8 p
صباح الخير يا أساتذة. إزيّكم؟

%%K t05-01-9 p
\textsc{المهندس روفو}. --- كويسين، الحمد لله. صباح الخير يا يوحنا، ما أكبر
العيّل ده!

%%K t05-01-10 p
\textsc{العم}. --- أيوه والله، بقى راجل صغير. وهو جدع أوي، لازم تفهّمه كل حاجة
بيشوفها. والنهارده عايز يعرف البيت بيتبني إزاي.

%%K t05-tit-2 sub
\VUpk{10.2pt}{40}{العمال.}
\VUsaut{3.0mm}

%%K t05-01-11 p
\textsc{الأسطى بلانكو}. --- طيب، تعالى معايا يا حبيبي، هشرح لك حالًا، علشان
أنا بحبّ أوي العيال اللي عايزة تتعلّم.

%%K t05-01-12 p
شايف العامل ده اللي ماسك \VUgras{كوريك} \textsuperscript{(82)} في إيده: ده
\VUgras{الحفّار} \textsuperscript{(81)}. بيحفر \VUgras{خندق} \textsuperscript{(46)} علشان يحطّ
مواسير المية. وهو كمان حفر الأرض في مكان البيت، علشان البنّا يقيم
\VUgras{الحيطان} الأربعة. والجزء من الحيطان اللي في الأرض اسمه
\VUgras{الأساسات} \textsuperscript{(2)}. والفراغ اللي بين الأساسات بيعمل
\VUgras{البدروم} \textsuperscript{(3)}. والبدروم ده عنده شباك صغير اسمه \VUgras{شباك
البدروم} \textsuperscript{(5)}، و\VUgras{باب} \textsuperscript{(4)}، وسلّم.

%%K t05-01-13 p
و\VUgras{البنّا} \textsuperscript{(108)} كمّل بناء الحيطان وساب فتحات لـ\VUgras{الأبواب}
\textsuperscript{(8)} ولـ\VUgras{الشبابيك} \textsuperscript{(17)}.

%%K t05-01-14 p
\textsc{يوحنا}. --- بس أنا مش شايف البنّا ده!

%%K t05-01-15 p
\textsc{الأسطى بلانكو}. --- هو خلّص شغله في البيت دلوقت. هو جنب
\VUgras{الإسطبل} \textsuperscript{(54)}، على \VUgras{السقالة} \textsuperscript{(110)} بتاعته،
بيبني حيطة \VUgras{المغسلة} \textsuperscript{(56)}.

%%K t05-01-16 p
وعنده مسطرين، وطشت، و\VUgras{شاقول} \textsuperscript{(109)}. بس إنت مش هتفتكر الأسامي دي.

%%K t05-01-17 p
\textsc{يوحنا}. --- أكيد هفتكرها، علشان أنا حالًا هطلب من عمّي مسطرين صغير
وطشت، وأنا بنفسي هعمل شاقول بخيط.

%%K t05-tit-3 sub
\VUsaut{3.0mm}
\VUpk{10.2pt}{40}{المواد.}
\VUsaut{3.0mm}

%%K t05-01-18 p
\textsc{العم}. --- آه! جميل أوي. بس قول لي يا يوحنا، إيه اللي لازم علشان يتبني
بيت؟

%%K t05-01-19 p
\textsc{يوحنا}. --- لازم \VUgras{حجر منحوت} \textsuperscript{(59)} و\VUgras{مونة}
\textsuperscript{(65)}.

%%K t05-01-20 p
\textsc{العم}. --- بالظبط، بُصّ الراجل ده، بالبرنيطة الكبيرة، على
\VUgras{الكارّو} \textsuperscript{(136)} بتاعه. ده \VUgras{العربجي} \textsuperscript{(135)}.
كل يوم
بيجيب هنا \VUgras{كتل حجر} \textsuperscript{(60)}. وبينزّل حمولته في الحوش،
و\VUgras{النحّاتين} \textsuperscript{(83)} بينحتوا الكتل ويقطّعوها بـ\VUgras{منشار الحجر}
\textsuperscript{(85)}. و\VUgras{الشيّال} \textsuperscript{(146)} بيوديها للبنّا اللي
بيركّبها.

%%K t05-01-21 p
وفي نفس الوقت، \VUgras{بتاع المونة} \textsuperscript{(87)} بيحضّر المونة. وهو محتاج
\VUgras{رملة} \textsuperscript{(62)}، و\VUgras{جير} \textsuperscript{(63)}، و\VUgras{مية}
\textsuperscript{(64)}. والجير جنبه في \VUgras{شكارة} \textsuperscript{(71)}، و\VUgras{صبي
البنّا} \textsuperscript{(111)} بيجيب له الرملة على \VUgras{عربية اليد}
\textsuperscript{(112)}، و\VUgras{الصبي} \textsuperscript{(113)} واقف على الطرمبة
\textsuperscript{(58)}. وبيملا \VUgras{جردل} \textsuperscript{(114)}. والمونة هتبقى جاهزة
حالًا. و\VUgras{شيّال المونة} \textsuperscript{(107)} بياخدها ويطلّعها فوق، على السلّم،
على السقالة، اللي عليها بنّا بيخلّص الناحية دي من البيت.

%%K t05-01-22 p
ودلوقت، اسم العامل اللي بيشتغل في \VUgras{السقف} \textsuperscript{(31)} إيه؟

%%K t05-01-23 p
\textsc{يوحنا}. --- ده \VUgras{بتاع السقف} \textsuperscript{(118)}.

%%K t05-tit-4 sub
\VUsaut{3.0mm}
\VUpk{10.2pt}{40}{سقف البيت.}
\VUsaut{3.0mm}

%%K t05-01-24 p
\textsc{الأسطى بلانكو}. --- أيوه، بس الأول بييجي \VUgras{النجّار}
\textsuperscript{(115)}.

%%K t05-01-25 p
النجّار بيشتغل \VUgras{تسقيفة الخشب} \textsuperscript{(35)}. وبيركّب \VUgras{الكمرات}
\textsuperscript{(37)} بـ\VUgras{مسامير خشب} \textsuperscript{(117)}. والعمال دول، فوق هناك،
ببنطلوناتهم الواسعة القطيفة، دول نجّارين. واحد ماسك كمرة صغيرة، والتاني
بيقطّع مسمار بـ\VUgras{بلطته} \textsuperscript{(116)}. واللي جنب \VUgras{المدخنة}
\textsuperscript{(41)} ده بتاع السقف. وهو بيسمّر

%%K t05-noto-1 noto
\VUnotes{143.2mm}{%
(*) \VUgras{Balk-o} = D. \textit{Balken}، E. \textit{joist}، F. \textit{solive}،
I. \textit{travicello}، R. \textit{balka}، S. Port. \textit{viga}. جذر فنّي لسه
ما اتاخدش لغاية دلوقت، بس نقترحه علشان نترجم معنى الكلمة الفرنساوية
\textit{solive}، اللي مش هي trabo أي F. \textit{poutre}، ولا trabeto. دي عصاية
خشب مربّعة بتشيل أرضية الخشب والسقف.%
}

%%K t05-01-25 p suite
الشرايح على (*) الكمرات الصغيرة، و\VUgras{الأردواز} \textsuperscript{(66)} على الشرايح.
وساعات بيحطّ قرميد.

%%K t05-01-26 p
سقف الإسطبل \VUgras{بالقرميد} \textsuperscript{(55)}، وسقف البيت \VUgras{بالأردواز}
\textsuperscript{(32)}، وسقف الفيرندة \VUgras{بالزنك} \textsuperscript{(24)}.

%%K t05-01-27 p
\textsc{يوحنا}. --- وبتاع السقف بيشتغل السقوف الزنك كمان؟

%%K t05-01-28 p
\textsc{الأسطى بلانكو}. --- لأ، ده \VUgras{السمكري} \textsuperscript{(121)} أو
\VUgras{بتاع الرصاص}، اللي بيحطّ \VUgras{الطاقة} \textsuperscript{(38)} على سقف البيت.
وهو كمان بيحطّ \VUgras{مزاريب المطر} \textsuperscript{(43)} و\VUgras{مزاريب السقف}
\textsuperscript{(42)}. وهو بيلحّم الزنك والصفيح. وهو اللي ثبّت \VUgras{مانعة
الصواعق} \textsuperscript{(40)} و\VUgras{ديك الريح} \textsuperscript{(39)} على
\VUgras{الجملون}
\textsuperscript{(36)}.

%%K t05-01-29 p
\textsc{يوحنا}. --- يبقى البيت خلص كده؟

%%K t05-tit-5 sub
\VUsaut{3.0mm}
\VUpk{10.2pt}{40}{العُدّة.}
\VUsaut{3.0mm}

%%K t05-01-30 p
\textsc{الأسطى بلانكو}. --- مش خالص. \VUgras{المبيّض} \textsuperscript{(99)} بيبني
بـ\VUgras{الطوب} \textsuperscript{(61)} الحيطان اللي بتقسّمه لأُوَض مختلفة. وبيدهن
بـ\VUgras{الجبس} \textsuperscript{(70)} \VUgras{الأسقف} \textsuperscript{(26)} والحيطان.
ودلوقت
بيشتغل سقف \VUgras{الفيرندة} \textsuperscript{(33)}. وعنده، زي البنّا،
\VUgras{مسطرين} \textsuperscript{(100)} و\VUgras{طشت} \textsuperscript{(101)}.

%%K t05-01-31 p
وبعد كده النجّار بيشتغل الأبواب، والشبابيك، و\VUgras{الباركيه}
\textsuperscript{(16)}، و\VUgras{الشيش} \textsuperscript{(19)}.

%%K t05-01-32 p
ودلوقت هو بيشتغل في الحوش على \VUgras{بنك الشغل} \textsuperscript{(90)} بتاعه. وعنده
\VUgras{منشار} \textsuperscript{(94)} علشان يقطّع الخشب \textsuperscript{(69)}،
و\VUgras{فارة} \textsuperscript{(91)}، و\VUgras{ملزمة} \textsuperscript{(92)}، و\VUgras{أزميل}
\textsuperscript{(94 bis)}، و\VUgras{فرجار} \textsuperscript{(95)}، و\VUgras{شاكوش}
\textsuperscript{(95 bis)} خشب.

%%K t05-01-33 p
\textsc{يوحنا}. --- آه! بس النجارة أحلى من البناء. يا عمّي، هتشتري لي بنك شغل،
ومنشار، وفارة، علشان أنا قريّب هعمل عشّة لميدور.

%%K t05-01-34 p
\textsc{العم}. --- ماشي يا حبيبي.

%%K t05-01-35 p
\textsc{يوحنا}. --- ما أسعدني! ومين اللي واقف جنب باب الدخول؟

%%K t05-tit-6 sub
\VUsaut{3.0mm}
\VUpk{10.2pt}{40}{الدهان.}
\VUsaut{3.0mm}

%%K t05-01-36 p
\textsc{الأسطى بلانكو}. --- ده \VUgras{النقّاش} \textsuperscript{(102)}. واقف على سلّمه
وعنده علبة \VUgras{بوية} \textsuperscript{(103)} و\VUgras{فرشته} \textsuperscript{(104)}.
وبيدهن
خشب باب الدخول علشان يفضل مدة أطول.

%%K t05-01-37 p
وهو كمان بيلزق الورق في الشقق.

%%K t05-01-38 p
و\VUgras{المنجّد} \textsuperscript{(107)} اللي على \VUgras{السلّم} \textsuperscript{(106)}،
على
\VUgras{البلكونة} \textsuperscript{(28)}، بيركّب \VUgras{ستارة} \textsuperscript{(29)}.

%%K t05-01-39 p
والعامل الواقف جنب الشباك الكبير هو \VUgras{القزّاز} \textsuperscript{(96)}. بيثبّت
\VUgras{ألواح القزاز} \textsuperscript{(18)} بـ\VUgras{المعجون} \textsuperscript{(73)}.
والعامل
التاني، الراكع جنب باب البدروم، هو \VUgras{بتاع الكوالين} \textsuperscript{(97)}. بيركّب
\VUgras{كالون} \textsuperscript{(6)}. و\VUgras{مبرده} \textsuperscript{(98)} جنبه.

%%K t05-01-40 p
\textsc{يوحنا}. --- والعامل اللي بيخبط بـ\VUgras{شاكوشه} \textsuperscript{(84)} على حديد،
ده كمان بتاع كوالين؟

%%K t05-01-41 p
\textsc{الأسطى بلانكو}. --- بالأصحّ، ده حدّاد \textsuperscript{(125)}. بيشكّل
\VUgras{حدايد} \textsuperscript{(67)} لـ\VUgras{الكهربجي} \textsuperscript{(124)} اللي على سقف
\VUgras{الجراج} \textsuperscript{(51)}. وعنده \VUgras{كور} \textsuperscript{(126)}،
و\VUgras{سندال} \textsuperscript{(127)}، و\VUgras{زرديّة} \textsuperscript{(128)}.

%%K t05-01-42 p
والكهربجي بيثبّت \VUgras{سلك} \VUgras{النحاس} \textsuperscript{(44)} بتاع التليفون.

%%K t05-tit-7 sub
\VUpk{10.2pt}{40}{الجنينة.}
\VUsaut{3.0mm}

%%K t05-01-43 p
\textsc{العم}. --- في آخر \VUgras{الحوش} \textsuperscript{(45)}، جنب \VUgras{السور
الحديد} \textsuperscript{(47)} و\VUgras{البوابة} \textsuperscript{(48)}، إنت شايف
\VUgras{الجنايني} \textsuperscript{(129)}. بيحضّر \VUgras{الجنينة الصغيرة}
\textsuperscript{(49)}. قلّب الأرض بـ\VUgras{فاسه} \textsuperscript{(131)}، وبيبذر البذور
ويمشّط بـ\VUgras{المشط} \textsuperscript{(130)}. وبعد كده، بـ\VUgras{الرشّاشة}
\textsuperscript{(132)} بتاعته، هيسقي \VUgras{الأحواض} \textsuperscript{(150)} والنباتات
هتكبر.

%%K t05-01-44 p
و\VUgras{السايس} \textsuperscript{(133)} اللي حالًا نضّف الحصان وأكّله، بينضّف
\VUgras{العربية} \textsuperscript{(52)} بسفنجة، و\VUgras{طرمبة إيد} \textsuperscript{(134)}،
وجردل مية. وبعد كده هيدخّلها الجراج ويقفل الباب بـ\VUgras{الترباس}
\textsuperscript{(53)} التخين.

%%K t05-01-45 p
أهو إنت مبسوط، أنا متصوّر، أخدت شرح كفاية.

%%K t05-01-46 p
\textsc{يوحنا}. --- أيوه يا عمّي، بس نفسي أشوف البيت من جوّه.

%%K t05-01-47 p
\textsc{العم}. --- ده هيبقى بكرة. المهندس روفو هيشرح لك الرسم ويقول لك اسم كل
أوضة.

%%K t05-tit-8 sub
\VUsaut{3.0mm}
\VUpk{10.2pt}{40}{ترتيب البيت من جوّه.}
\VUsaut{2.4mm}

%%K t05-apar-2 apar
{\centering\textit{(بُصّ الرسم.)}\par}
\VUsaut{2.4mm}

%%K t05-01-48 p
\textsc{المهندس}. --- تعالى يا يوحنا، حالًا هوريك أجزاء البيت المختلفة.

%%K t05-01-49 p
تعالى ندخل \VUgras{الدور الأرضي} \textsuperscript{(7)}. أهو زرار \VUgras{الجرس}
\textsuperscript{(11)}. بنطلع تلات \VUgras{درجات} \textsuperscript{(14)} بتاعة \VUgras{السلّم
الخارجي} \textsuperscript{(9)}، وبنعدّي \VUgras{عتبة} \textsuperscript{(10)} \VUgras{باب
الدخول}
\textsuperscript{(8)}، وأهو إحنا عند رجل \VUgras{السلّم} \textsuperscript{(13)}.

%%K t05-01-50 p
\textsc{يوحنا}. --- دي الطرقة.

%%K t05-01-51 p
\textsc{المهندس}. --- لأ، ده \VUgras{المدخل} \textsuperscript{(12)}.
و\VUgras{الطرقة} \textsuperscript{(a)} في الآخر. على اليمين، أهو \VUgras{الصالون}
\textsuperscript{(b)} و\VUgras{أوضة السفرة} \textsuperscript{(c)}؛ وقصاد أوضة السفرة
\VUgras{المطبخ} \textsuperscript{(d)} و\VUgras{أوضة الخدم} \textsuperscript{(e)}؛ وبعدين قدّام
\VUgras{المكتب} \textsuperscript{(f)}. و\VUgras{الفيرندة} \textsuperscript{(23)}
بـ\VUgras{بابها القزاز} \textsuperscript{(25)} جنب الصالون.

%%K t05-01-52 p
حالًا هنطلع السلّم. اتمسّك في \VUgras{الدرابزين} \textsuperscript{(15)} وعُدّ الدرجات.
هما أربعتاشر. أهي \VUgras{البسطة} \textsuperscript{(m)}، إحنا على \VUgras{الدور الأول}
\textsuperscript{(27)}.

%%K t05-tit-9 sub
\VUsaut{3.0mm}
\VUpk{10.2pt}{40}{الدور الأول.}
\VUsaut{3.0mm}

%%K t05-01-53 p
وقصاد السلّم، في \VUgras{بيت الأدب} \textsuperscript{(20)} و\VUgras{أوضة التواليت}
\textsuperscript{(j)}. وعلى اليمين \VUgras{أوضة نوم} \textsuperscript{(g)} حلوة بشبّاكين على
\VUgras{واجهة} \textsuperscript{(1)} البيت. وعلى الشمال، في \VUgras{الحمّام}
\textsuperscript{(1)} و\VUgras{أوضة الضيوف} \textsuperscript{(i)} بـ\VUgras{ألكوف}
\textsuperscript{(h)} كبير. وجنب أوضة النوم في \VUgras{أوضة العيال}
\textsuperscript{(k)}. وهي على \VUgras{تراس} \textsuperscript{(21)} واسع. وتحت التراس ده في
بيت أدب تاني \textsuperscript{(20)}، وعلى الحيطة، أهو \VUgras{الجرس}
\textsuperscript{(22)} اللي بيدق للعشا.

%%K t05-01-54 p
\textsc{يوحنا}. --- تعالى نطلع \VUgras{الدور التاني}.

%%K t05-01-55 p
\textsc{المهندس}. --- مفيش دور تاني. تحت السقف مافيش غير \VUgras{أُوَض تحت
السطوح} \textsuperscript{(33)} والمخزن \textsuperscript{(34)}. خلّصنا الجولة، نقدر
ننزل تاني.

%%K t05-01-56 p
\textsc{يوحنا}. --- شكرًا يا أستاذ، ده مشوّق أوي. حالًا هحكي لأبويا كل اللي
شفته.

\VUsaut{4.0mm}
\VUfilet{22mm}
